\chapter{Background}
\label{chap:background}

This chapter provides the background needed to follow the rest of the
thesis. We cover finite fields, elliptic curves, bilinear pairings,
and give a brief overview of each of the four schemes. The detailed
mathematics and pseudocode for each scheme are presented in their
respective chapters.
 
% ────────────────────────────────────────────────────────────
\section{Finite Fields and Elliptic Curves}
\label{sec:fields-curves}
 
A \emph{prime field} $\mathbb{F}_p$ is the set of integers
$\{0, 1, \ldots, p-1\}$ where $p$ is a prime, with addition and
multiplication done modulo $p$. Elliptic curve cryptography works
over points on a curve defined over such a field.
 
An elliptic curve over $\mathbb{F}_p$ in short Weierstrass form is:
\begin{equation}
    E : y^2 = x^3 + ax + b,
    \quad a, b \in \mathbb{F}_p,
    \quad 4a^3 + 27b^2 \neq 0.
    \label{eq:weierstrass}
\end{equation}
 
The points on this curve form an abelian group under the standard
point addition law. The core operation in all protocols studied here
is \emph{scalar multiplication}: given a point $P$ and an integer
$k$, computing $[k]P = P + P + \cdots + P$ ($k$ times). The cost of
this operation directly determines how fast commitment, proving, and
verification are in practice.
 
% ────────────────────────────────────────────────────────────
\section{Bilinear Pairings}
\label{sec:pairings}
 
Let $\mathbb{G}_1$, $\mathbb{G}_2$ be two groups of prime order $r$
on an elliptic curve, and $\mathbb{G}_T$ a multiplicative group of
the same order. A bilinear pairing is a map:
\begin{equation}
    e : \mathbb{G}_1 \times \mathbb{G}_2 \longrightarrow \mathbb{G}_T
\end{equation}
that satisfies $e([a]P,\, [b]Q) = e(P,Q)^{ab}$ for all scalars
$a, b$. This bilinearity property is what allows pairing-based
schemes like KZG, MKZG, and Dory to verify polynomial identities
using only committed values, without the verifier ever seeing the
polynomial. Pairing computations are significantly more expensive
than scalar multiplications, so the number of pairings required
during verification is a key metric for comparing schemes.
 
% ────────────────────────────────────────────────────────────
\section{Elliptic Curves Used in This Work}
\label{sec:curves}
 
All three curves used in this thesis are pairing-friendly curves with
embedding degree 12, meaning they support efficient pairing
computations. Table~\ref{tab:curves} summarises their key parameters.
 
\begin{table}[h!]
\centering
\caption{Parameters of the three elliptic curves used in this thesis.}
\label{tab:curves}
\begin{tabular}{lccc}
\toprule
\textbf{Property} & \textbf{BLS12-381} & \textbf{BN254} & \textbf{BLS12-377} \\
\midrule
Field size (bits)     & 381  & 254      & 377  \\
Subgroup order (bits) & 255  & 254      & 253  \\
Security level (bits) & 128  & 100--110 & 128  \\
Trusted setup needed  & Yes  & Yes      & Yes  \\
Designed for          & Zcash & Ethereum & Recursion \\
\bottomrule
\end{tabular}
\end{table}
 
\textbf{BLS12-381} was designed by Sean Bowe in 2017 for Zcash. It
offers 128-bit security and is widely used in production SNARK
systems.
 
\textbf{BN254} was introduced by Barreto and Naehrig in 2005 and was
the default curve in early Ethereum zero-knowledge applications. Its
security is slightly lower than the other two due to the
Kim-Barbulescu attacks on the target field, but its smaller field
size makes arithmetic faster in software.
 
\textbf{BLS12-377} was designed specifically to support recursive
proof composition, where one proof verifies another proof. Its
subgroup order matches the base field of the outer curve BW6-761,
which makes recursive verification efficient.
 
% ────────────────────────────────────────────────────────────
\section{Polynomial Commitment Schemes}
\label{sec:pcs-overview}
 
A polynomial commitment scheme allows a prover to commit to a
polynomial and later prove its evaluation at any point chosen by
the verifier. Formally it consists of four algorithms:
\textsc{Setup}, \textsc{Commit}, \textsc{Open}, and \textsc{Verify}.
A scheme must be correct, binding, and optionally hiding.
 
The four schemes studied in this thesis differ in whether they need
a trusted setup, how large their proofs are, and how expensive
proving and verification are. Table~\ref{tab:scheme-comparison}
gives a high-level comparison.
 
\begin{table}[h!]
\centering
\caption{High-level comparison of the four polynomial commitment
schemes. $d$ is the polynomial degree and $\mu$ is the number of
variables for multilinear polynomials.}
\label{tab:scheme-comparison}
\resizebox{\textwidth}{!}{%
\begin{tabular}{lccccc}
\toprule
\textbf{Scheme} & \textbf{Trusted Setup} & \textbf{Proof Size}
               & \textbf{Prover Cost} & \textbf{Verifier Cost}
               & \textbf{Uses Pairings} \\
\midrule
KZG             & Yes, $O(d)$        & $O(1)$        & $O(d)$
                & $O(1)$             & Yes \\
Multilinear KZG & Yes, $O(2^\mu)$    & $O(\mu)$      & $O(2^\mu)$
                & $O(\mu)$           & Yes \\
Bulletproofs    & No                 & $O(\log d)$   & $O(d)$
                & $O(d)$             & No  \\
Dory            & No                 & $O(\log d)$   & $O(d)$
                & $O(\log d)$        & Yes \\
\bottomrule
\end{tabular}}
\end{table}
 
\textbf{KZG} \cite{kate2010} is the most widely deployed scheme in
practice. It requires a trusted setup but produces constant-size
proofs and constant-time verification, making it very attractive for
applications where verification speed matters.
 
\textbf{Multilinear KZG} \cite{papamanthou2013} extends KZG to
polynomials over the Boolean hypercube. It is used in systems like
HyperPlonk \cite{hyperplonk2022} that work with multilinear
arithmetic.
 
\textbf{Bulletproofs} \cite{bulletproofs2018} require no trusted
setup at all. The trade-off is that verification cost grows linearly
with the polynomial degree, which makes them slower to verify than
pairing-based schemes at large sizes.
 
\textbf{Dory} \cite{dory2021} also requires no trusted setup and
achieves $O(\log d)$ verification using pairings, sitting between
Bulletproofs and KZG on the transparency-efficiency spectrum.
 
Each scheme is described in full detail, with pseudocode and
benchmark results, in its own chapter.